% Experiment setup. Drops into the main paper with % Experiment setup. Drops into the main paper with % Experiment setup. Drops into the main paper with % Experiment setup. Drops into the main paper with \input{experiment_setup};
% also compiles standalone (the guard falls back to article when the workshop
% style is not on TEXINPUTS).
\documentclass{article}

\IfFileExists{neurips_2026.sty}{\usepackage[dblblindworkshop]{neurips_2026}}{%
  % standalone fallback: lmodern gives microtype a scalable font to expand
  \usepackage[margin=1in]{geometry}\usepackage{lmodern}}

\usepackage[utf8]{inputenc}
\usepackage[T1]{fontenc}
\usepackage{hyperref}
\usepackage{url}
\usepackage{booktabs}
\usepackage{amsfonts}
\usepackage{amsmath}
\usepackage{amssymb}
\usepackage{microtype}
\usepackage{graphicx}
\usepackage{xcolor}

\title{Branchpoint: Towards Structured Multiverse Analysis}

\author{%
  Anonymous Author(s) \\
  Affiliation \\
  \texttt{email} \\
}

\begin{document}
\maketitle

\section{Experiment setup}
\label{sec:setup}

This section covers everything outside the sampler that the results depend on:
the datasets, what the analyst is asked, how the model is called, the arms, and
what happens to a result once it exists. The sampler itself is described in
Section~\ref{sec:setup} of the main text; it detects contestable decisions in a
recorded analysis, forks them under a neutral intervention, and grows a ban
list. It makes no validity judgement and never sees a standardization rule.

A third dataset (MARP, religiosity and well-being \cite{hoogeveen2023many})
remains configured but is \emph{dropped throughout}: not run, not standardized,
not plotted, not reported. Nothing below applies to it.

\subsection{Datasets}
\label{sec:setup-data}

Two ecology datasets, each with a published many-analyst comparison, both from
ManyEcoEvo \cite{gould2025same}.

\begin{table}[h]
\centering
\small
\begin{tabular}{lll}
\toprule
 & Blue tit (\texttt{bt}) & Eucalyptus (\texttt{eu}) \\
\midrule
Dimensions            & $3{,}720 \times 42$ & $351 \times 43$ \\
Observation unit      & one chick & one quadrat in one round \\
Clustering            & chicks within ${\sim}330$ rearing nests & quadrats within 18 properties \\
Human analyses coded  & 174 & 128 \\
Usable for $Z_r$      & 111 & 85 \\
\bottomrule
\end{tabular}
\caption{The two datasets. Blue-tit chicks are nested within rearing broods;
Eucalyptus quadrats are not nested in any comparable way, and physical quadrats
do \emph{not} recur across survey rounds. That difference propagates all the way
to the degrees-of-freedom rule in Section~\ref{sec:setup-layer2}.}
\label{tab:datasets}
\end{table}

Analysts answer the published question verbatim: \emph{``To what extent is the
growth of nestling blue tits (Cyanistes caeruleus) influenced by competition
with siblings?''} and \emph{``How does grass cover influence Eucalyptus spp.\
seedling recruitment?''}

\paragraph{Two files per dataset, doing different jobs.}
The analyst-facing data is what the model analyses. A separate comparator file
carries each human team's reported coefficient, standard error and adjusted
degrees of freedom, from which the crowd distribution is derived. Conflating
them is a live hazard, and the two are never interchangeable.

\paragraph{The Eucalyptus design description matters.}
An earlier snapshot of the question stated that fixed quadrats recur across
survey rounds and that the file has 43 columns. Neither is true: quadrat labels
are reused and are not persistent identifiers, and the supplied file has 38
usable columns. Analysts given the false description fit quadrat-clustered
repeated-measures models, which report an asymptotic Wald $z$ with no finite
error df and are then unstandardizable. Every run reported here uses the
corrected snapshot, hash-recorded per condition alongside the results.

\subsection{What the analyst is told}
\label{sec:setup-prompt}

One prompt per dataset, with a fixed structure: the published research question;
a summary of the study design; a column guide; the declared estimand in plain
words; and a frame policy stating what counts as in-frame. Then the working
instruction: analyse the data, report the fitted analysis's focal estimate and
uncertainty \emph{in its native model scale}, work in small checkable Python
scripts, and report a note describing exactly what model, variables and
exclusions were used.

\paragraph{The analyst is never told the standardization rule.}
No Fisher-$z$ formula, no $t$/df transformation map, no degrees-of-freedom
convention. This boundary is marked in code and structurally tested. It matters
because the transformation is exactly the kind of downstream convention that
would contaminate the analytic choices being measured.

This mirrors the benchmark. ManyEcoEvo analysts received the one-sentence
question and the data and nothing else; the information needed for
standardization was elicited \emph{afterwards}, in a separate follow-up survey
with published per-dataset instructions \cite{gould2025same}. Our pipeline
currently reconstructs that information from the analyst's free-text note rather
than asking for it, which is a limitation we return to in
Section~\ref{sec:setup-limits}.

\subsection{How the model is called}
\label{sec:setup-calls}

\paragraph{Loop.}
A files-as-state analysis loop gives the model two tools: \texttt{run\_code},
which writes and runs one Python analysis step in the project directory, and
\texttt{submit\_report}, which finishes with the focal estimate and standard
error. Every step appends to a durable transcript. The model builds its analysis
across small checkable steps rather than one script.

\paragraph{Provider interface.}
All arms use the provider-native interface. Per call we record a secret-free
request object, the complete SDK-serialized response, opaque reasoning blocks,
usage, and every host tool result. This is an SDK-level record, not an
HTTP-wire capture, and it does not expose hidden model state.

\paragraph{Models.}
Three analysts (\texttt{gpt-5.6-sol}, \texttt{claude-opus-5},
\texttt{glm-5.3}) and one judge (\texttt{gpt-5.6-sol}), all at low reasoning
effort. \texttt{glm-5.3} is pinned to a single provider for continuation so the
native history is stable across a fork. Cost is metered per call, tagged by
role, written to a ledger, and re-derived from token counts at one rate table; a
spend guard aborts a run above a ceiling. The full grid below cost \$105.

\subsection{Arms}
\label{sec:setup-arms}

Three arms, sharing the same loop, tools and question.

\begin{table}[h]
\centering
\small
\begin{tabular}{lp{9.4cm}}
\toprule
Arm & What it does \\
\midrule
Naive & $n$ independent analyses from the same prompt under different seeds. The
control: no decision detection, no forking, no intervention. \\
Branchpoint & Detects decisions in a completed trajectory, copies state after the
complete response, appends a neutral intervention, and continues from the exact
provider-native history. Branch points are drawn at random from the open set. \\
BpDynamic & Identical to Branchpoint except allocation: open decisions are
sorted by branches taken plus strikes, and the least-tried are drawn first, so
budget is spent on uncovered decisions. \\
\bottomrule
\end{tabular}
\caption{The three arms. Branchpoint and BpDynamic differ only in how branch
points are allocated, not in how a branch is produced.}
\label{tab:arms}
\end{table}

The intervention is one neutral sentence naming the decision and listing the
approaches already used: \emph{``For \textlangle decision\textrangle, do not use
any approach in this JSON list: [\ldots]. Choose a different defensible
approach; everything else remains yours to decide.''}

\paragraph{Grid as run.}
$\{\textrm{sol}, \textrm{opus}, \textrm{glm}\} \times \{\textrm{bt},
\textrm{eu}\} \times \{\textrm{Naive}, \textrm{Branchpoint}\}$, plus
$\textrm{sol} \times \{\textrm{bt}, \textrm{eu}\} \times \textrm{BpDynamic}$:
14 conditions, $n = 50$ each, ${\sim}700$ analyses.

\subsection{Layer 2: what happens to a result}
\label{sec:setup-layer2}

Layer 2 is host-side, runs only after Layer 1 is banked, and is the sole owner
of the conversion to a common scale. It never runs or extends an analysis agent.
Four stages, each digest-bound so an unchanged input is never re-judged.

\paragraph{(i) Extract.}
A judge reads what the run printed and returns a fixed schema: number of
observations, the unit, the design, exactly one of a Student $t$ or an
asymptotic Wald $z$ (never relabelled), the residual df, a separate finite
reference df used for inference such as clusters minus one, and any directly
reported standardized effect. The digest binds prompt, schema and judge model.

\paragraph{(ii) Standardize --- deterministic, no model call.}
The partial correlation and its transform are
\begin{equation}
r = \frac{t}{\sqrt{t^{2} + \mathit{df}}},
\qquad
Z_r = \operatorname{atanh}(r),
\qquad
\mathrm{SE}(Z_r) = \frac{1}{\sqrt{\mathit{df}}} .
\label{eq:standardize}
\end{equation}
A row reporting both $t$ and $z$ is blocked as ambiguous. A row with a statistic
but no usable df is blocked, but not immediately dropped: it falls through to a
directly reported effect, taken as-is when already on the Fisher-$z$ scale and
via $\operatorname{atanh}$ when $|{\cdot}| < 1$. Only when no usable direct
effect exists does the blocked reason become the row's terminal state, and the
row carries that exact reason rather than disappearing.

Note that $\mathrm{SE}(Z_r)$ depends only on the degrees of freedom: the
analyst's own reported standard error is deliberately \emph{not} used, following
the same convention as the benchmark \cite{gould2025same}. A consequence worth
stating plainly is that our intervals carry no information beyond df, so within
a condition they are nearly identical.

\paragraph{Degrees of freedom are per-dataset policy.}
Equation~\ref{eq:standardize} requires the df that the comparator itself used,
and the two datasets do not agree. Measured from each dataset's own human
panel, the median ratio of adjusted df to sample size is $0.126$ for blue tit
(90 of 96 analyses on a cluster scale) and $0.936$ for Eucalyptus (40 of 55 on a
row scale). They differ because chicks nest within broods and quadrats do not.

Policy therefore decides \emph{precedence}, not merely a fallback: on a
row-scale comparator the residual df wins, then $n - 2$, then any reference df;
on a cluster-scale comparator the reference df wins, then the residual df. This
matters because $\mathrm{SE} = 1/\sqrt{\mathit{df}}$, so a row on a cluster df
sitting beside a row on a row df receives an interval several-fold different for
the same analysis type, and the \emph{gap between the two scales} is then read
as spread. Before precedence was enforced, Eucalyptus widths were inflated
roughly threefold with no change in the number of rows retained.

Where the comparator is row-scale and a generalized linear model reports an
asymptotic Wald $z$ with no classical error df, $n - 2$ is derived. It is the
conservative end, since it maximizes df and so cannot inflate $|r|$, and the
number of parameters is immaterial at these sample sizes: at $n = 350$,
$z = -1.12$, $Z_r$ is $-0.0602$ at $\mathit{df} = 345$ against $-0.0615$ at
$\mathit{df} = 330$. The df \emph{family} is what matters --- a cluster df of 17
would give $-0.262$, over four times larger. Where the comparator is
cluster-scale no cluster count is ever guessed; a row carrying only a row-scale
df stays on it and is reported as such.

\paragraph{(iii) Scope.}
One judgement per standardizable row, against the declared estimand and frame
only. The prompt is explicit that a different covariate set, model family or
estimator is an analytical choice and \emph{not} a scope change, and that
neither is the aggregation level: individuals, clusters and group means are all
reasonable ways to answer one question, and the human panels contain all three.
In this run the screen removed almost nothing.

\paragraph{(iv) Finalize.}
The reportable spread over rows that are both standardizable and in scope. Every
excluded row is retained with its reason, and the counts are reported
separately, so attrition is always visible beside the width --- which matters,
because the two are not independent.

\paragraph{Manual overrides are not used here.}
The pipeline supports a per-condition override file that short-circuits
\emph{before} the deterministic path, supplying a value and method directly.
None of the 14 conditions reported here carries one; every value below is
produced by Equation~\ref{eq:standardize}. This is worth stating because a
condition carrying such a file has rows that no subsequent code fix can reach.

\subsection{Known limitations}
\label{sec:setup-limits}

\begin{itemize}
\item \textbf{Standardization information is reconstructed, not elicited.} The
benchmark asked its analysts for the statistic and df in a follow-up survey; we
parse a free-text note. Analyses whose author wrote ``$z$'' rather than ``$t$
with df'' were being dropped by reporting convention rather than by anything
analytic. The principled fix is a post-analysis elicitation stage.
\item \textbf{Coverage is capped, not converged.} With $n = 50$ against roughly
200 detected decisions, no decision is explored to exhaustion: nothing closes,
and every run stops at the cap. Reported widths are lower bounds from a partial
sample, and coverage belongs beside them.
\item \textbf{Ban lists are per-trajectory.} The same decision axis explored
from a different parent starts with an empty ban list, so distinct branches can
land on identical specifications.
\item \textbf{Blue-tit attrition is uneven.} One condition retains 10 of 50.
Rows reporting a cluster-robust $z$ with no stated cluster count remain
unstandardizable by design, rather than having a cluster count guessed for them.
\item \textbf{Our analyses are far more homogeneous than the humans'.} Within a
condition the standard-error coefficient of variation is $0.002$--$0.09$ against
the crowd's $0.58$--$0.65$, and roughly a third of properly-clustered human
analyses report $|t| < 4$ against about 2\% of ours. No sampler or
standardization change reaches this. It is a property of the analyst, and it is
the central finding rather than a defect.
\end{itemize}

\bibliographystyle{plain}
\bibliography{refs}

\end{document}
;
% also compiles standalone (the guard falls back to article when the workshop
% style is not on TEXINPUTS).
\documentclass{article}

\IfFileExists{neurips_2026.sty}{\usepackage[dblblindworkshop]{neurips_2026}}{%
  % standalone fallback: lmodern gives microtype a scalable font to expand
  \usepackage[margin=1in]{geometry}\usepackage{lmodern}}

\usepackage[utf8]{inputenc}
\usepackage[T1]{fontenc}
\usepackage{hyperref}
\usepackage{url}
\usepackage{booktabs}
\usepackage{amsfonts}
\usepackage{amsmath}
\usepackage{amssymb}
\usepackage{microtype}
\usepackage{graphicx}
\usepackage{xcolor}

\title{Branchpoint: Towards Structured Multiverse Analysis}

\author{%
  Anonymous Author(s) \\
  Affiliation \\
  \texttt{email} \\
}

\begin{document}
\maketitle

\section{Experiment setup}
\label{sec:setup}

This section covers everything outside the sampler that the results depend on:
the datasets, what the analyst is asked, how the model is called, the arms, and
what happens to a result once it exists. The sampler itself is described in
Section~\ref{sec:setup} of the main text; it detects contestable decisions in a
recorded analysis, forks them under a neutral intervention, and grows a ban
list. It makes no validity judgement and never sees a standardization rule.

A third dataset (MARP, religiosity and well-being \cite{hoogeveen2023many})
remains configured but is \emph{dropped throughout}: not run, not standardized,
not plotted, not reported. Nothing below applies to it.

\subsection{Datasets}
\label{sec:setup-data}

Two ecology datasets, each with a published many-analyst comparison, both from
ManyEcoEvo \cite{gould2025same}.

\begin{table}[h]
\centering
\small
\begin{tabular}{lll}
\toprule
 & Blue tit (\texttt{bt}) & Eucalyptus (\texttt{eu}) \\
\midrule
Dimensions            & $3{,}720 \times 42$ & $351 \times 43$ \\
Observation unit      & one chick & one quadrat in one round \\
Clustering            & chicks within ${\sim}330$ rearing nests & quadrats within 18 properties \\
Human analyses coded  & 174 & 128 \\
Usable for $Z_r$      & 111 & 85 \\
\bottomrule
\end{tabular}
\caption{The two datasets. Blue-tit chicks are nested within rearing broods;
Eucalyptus quadrats are not nested in any comparable way, and physical quadrats
do \emph{not} recur across survey rounds. That difference propagates all the way
to the degrees-of-freedom rule in Section~\ref{sec:setup-layer2}.}
\label{tab:datasets}
\end{table}

Analysts answer the published question verbatim: \emph{``To what extent is the
growth of nestling blue tits (Cyanistes caeruleus) influenced by competition
with siblings?''} and \emph{``How does grass cover influence Eucalyptus spp.\
seedling recruitment?''}

\paragraph{Two files per dataset, doing different jobs.}
The analyst-facing data is what the model analyses. A separate comparator file
carries each human team's reported coefficient, standard error and adjusted
degrees of freedom, from which the crowd distribution is derived. Conflating
them is a live hazard, and the two are never interchangeable.

\paragraph{The Eucalyptus design description matters.}
An earlier snapshot of the question stated that fixed quadrats recur across
survey rounds and that the file has 43 columns. Neither is true: quadrat labels
are reused and are not persistent identifiers, and the supplied file has 38
usable columns. Analysts given the false description fit quadrat-clustered
repeated-measures models, which report an asymptotic Wald $z$ with no finite
error df and are then unstandardizable. Every run reported here uses the
corrected snapshot, hash-recorded per condition alongside the results.

\subsection{What the analyst is told}
\label{sec:setup-prompt}

One prompt per dataset, with a fixed structure: the published research question;
a summary of the study design; a column guide; the declared estimand in plain
words; and a frame policy stating what counts as in-frame. Then the working
instruction: analyse the data, report the fitted analysis's focal estimate and
uncertainty \emph{in its native model scale}, work in small checkable Python
scripts, and report a note describing exactly what model, variables and
exclusions were used.

\paragraph{The analyst is never told the standardization rule.}
No Fisher-$z$ formula, no $t$/df transformation map, no degrees-of-freedom
convention. This boundary is marked in code and structurally tested. It matters
because the transformation is exactly the kind of downstream convention that
would contaminate the analytic choices being measured.

This mirrors the benchmark. ManyEcoEvo analysts received the one-sentence
question and the data and nothing else; the information needed for
standardization was elicited \emph{afterwards}, in a separate follow-up survey
with published per-dataset instructions \cite{gould2025same}. Our pipeline
currently reconstructs that information from the analyst's free-text note rather
than asking for it, which is a limitation we return to in
Section~\ref{sec:setup-limits}.

\subsection{How the model is called}
\label{sec:setup-calls}

\paragraph{Loop.}
A files-as-state analysis loop gives the model two tools: \texttt{run\_code},
which writes and runs one Python analysis step in the project directory, and
\texttt{submit\_report}, which finishes with the focal estimate and standard
error. Every step appends to a durable transcript. The model builds its analysis
across small checkable steps rather than one script.

\paragraph{Provider interface.}
All arms use the provider-native interface. Per call we record a secret-free
request object, the complete SDK-serialized response, opaque reasoning blocks,
usage, and every host tool result. This is an SDK-level record, not an
HTTP-wire capture, and it does not expose hidden model state.

\paragraph{Models.}
Three analysts (\texttt{gpt-5.6-sol}, \texttt{claude-opus-5},
\texttt{glm-5.3}) and one judge (\texttt{gpt-5.6-sol}), all at low reasoning
effort. \texttt{glm-5.3} is pinned to a single provider for continuation so the
native history is stable across a fork. Cost is metered per call, tagged by
role, written to a ledger, and re-derived from token counts at one rate table; a
spend guard aborts a run above a ceiling. The full grid below cost \$105.

\subsection{Arms}
\label{sec:setup-arms}

Three arms, sharing the same loop, tools and question.

\begin{table}[h]
\centering
\small
\begin{tabular}{lp{9.4cm}}
\toprule
Arm & What it does \\
\midrule
Naive & $n$ independent analyses from the same prompt under different seeds. The
control: no decision detection, no forking, no intervention. \\
Branchpoint & Detects decisions in a completed trajectory, copies state after the
complete response, appends a neutral intervention, and continues from the exact
provider-native history. Branch points are drawn at random from the open set. \\
BpDynamic & Identical to Branchpoint except allocation: open decisions are
sorted by branches taken plus strikes, and the least-tried are drawn first, so
budget is spent on uncovered decisions. \\
\bottomrule
\end{tabular}
\caption{The three arms. Branchpoint and BpDynamic differ only in how branch
points are allocated, not in how a branch is produced.}
\label{tab:arms}
\end{table}

The intervention is one neutral sentence naming the decision and listing the
approaches already used: \emph{``For \textlangle decision\textrangle, do not use
any approach in this JSON list: [\ldots]. Choose a different defensible
approach; everything else remains yours to decide.''}

\paragraph{Grid as run.}
$\{\textrm{sol}, \textrm{opus}, \textrm{glm}\} \times \{\textrm{bt},
\textrm{eu}\} \times \{\textrm{Naive}, \textrm{Branchpoint}\}$, plus
$\textrm{sol} \times \{\textrm{bt}, \textrm{eu}\} \times \textrm{BpDynamic}$:
14 conditions, $n = 50$ each, ${\sim}700$ analyses.

\subsection{Layer 2: what happens to a result}
\label{sec:setup-layer2}

Layer 2 is host-side, runs only after Layer 1 is banked, and is the sole owner
of the conversion to a common scale. It never runs or extends an analysis agent.
Four stages, each digest-bound so an unchanged input is never re-judged.

\paragraph{(i) Extract.}
A judge reads what the run printed and returns a fixed schema: number of
observations, the unit, the design, exactly one of a Student $t$ or an
asymptotic Wald $z$ (never relabelled), the residual df, a separate finite
reference df used for inference such as clusters minus one, and any directly
reported standardized effect. The digest binds prompt, schema and judge model.

\paragraph{(ii) Standardize --- deterministic, no model call.}
The partial correlation and its transform are
\begin{equation}
r = \frac{t}{\sqrt{t^{2} + \mathit{df}}},
\qquad
Z_r = \operatorname{atanh}(r),
\qquad
\mathrm{SE}(Z_r) = \frac{1}{\sqrt{\mathit{df}}} .
\label{eq:standardize}
\end{equation}
A row reporting both $t$ and $z$ is blocked as ambiguous. A row with a statistic
but no usable df is blocked, but not immediately dropped: it falls through to a
directly reported effect, taken as-is when already on the Fisher-$z$ scale and
via $\operatorname{atanh}$ when $|{\cdot}| < 1$. Only when no usable direct
effect exists does the blocked reason become the row's terminal state, and the
row carries that exact reason rather than disappearing.

Note that $\mathrm{SE}(Z_r)$ depends only on the degrees of freedom: the
analyst's own reported standard error is deliberately \emph{not} used, following
the same convention as the benchmark \cite{gould2025same}. A consequence worth
stating plainly is that our intervals carry no information beyond df, so within
a condition they are nearly identical.

\paragraph{Degrees of freedom are per-dataset policy.}
Equation~\ref{eq:standardize} requires the df that the comparator itself used,
and the two datasets do not agree. Measured from each dataset's own human
panel, the median ratio of adjusted df to sample size is $0.126$ for blue tit
(90 of 96 analyses on a cluster scale) and $0.936$ for Eucalyptus (40 of 55 on a
row scale). They differ because chicks nest within broods and quadrats do not.

Policy therefore decides \emph{precedence}, not merely a fallback: on a
row-scale comparator the residual df wins, then $n - 2$, then any reference df;
on a cluster-scale comparator the reference df wins, then the residual df. This
matters because $\mathrm{SE} = 1/\sqrt{\mathit{df}}$, so a row on a cluster df
sitting beside a row on a row df receives an interval several-fold different for
the same analysis type, and the \emph{gap between the two scales} is then read
as spread. Before precedence was enforced, Eucalyptus widths were inflated
roughly threefold with no change in the number of rows retained.

Where the comparator is row-scale and a generalized linear model reports an
asymptotic Wald $z$ with no classical error df, $n - 2$ is derived. It is the
conservative end, since it maximizes df and so cannot inflate $|r|$, and the
number of parameters is immaterial at these sample sizes: at $n = 350$,
$z = -1.12$, $Z_r$ is $-0.0602$ at $\mathit{df} = 345$ against $-0.0615$ at
$\mathit{df} = 330$. The df \emph{family} is what matters --- a cluster df of 17
would give $-0.262$, over four times larger. Where the comparator is
cluster-scale no cluster count is ever guessed; a row carrying only a row-scale
df stays on it and is reported as such.

\paragraph{(iii) Scope.}
One judgement per standardizable row, against the declared estimand and frame
only. The prompt is explicit that a different covariate set, model family or
estimator is an analytical choice and \emph{not} a scope change, and that
neither is the aggregation level: individuals, clusters and group means are all
reasonable ways to answer one question, and the human panels contain all three.
In this run the screen removed almost nothing.

\paragraph{(iv) Finalize.}
The reportable spread over rows that are both standardizable and in scope. Every
excluded row is retained with its reason, and the counts are reported
separately, so attrition is always visible beside the width --- which matters,
because the two are not independent.

\paragraph{Manual overrides are not used here.}
The pipeline supports a per-condition override file that short-circuits
\emph{before} the deterministic path, supplying a value and method directly.
None of the 14 conditions reported here carries one; every value below is
produced by Equation~\ref{eq:standardize}. This is worth stating because a
condition carrying such a file has rows that no subsequent code fix can reach.

\subsection{Known limitations}
\label{sec:setup-limits}

\begin{itemize}
\item \textbf{Standardization information is reconstructed, not elicited.} The
benchmark asked its analysts for the statistic and df in a follow-up survey; we
parse a free-text note. Analyses whose author wrote ``$z$'' rather than ``$t$
with df'' were being dropped by reporting convention rather than by anything
analytic. The principled fix is a post-analysis elicitation stage.
\item \textbf{Coverage is capped, not converged.} With $n = 50$ against roughly
200 detected decisions, no decision is explored to exhaustion: nothing closes,
and every run stops at the cap. Reported widths are lower bounds from a partial
sample, and coverage belongs beside them.
\item \textbf{Ban lists are per-trajectory.} The same decision axis explored
from a different parent starts with an empty ban list, so distinct branches can
land on identical specifications.
\item \textbf{Blue-tit attrition is uneven.} One condition retains 10 of 50.
Rows reporting a cluster-robust $z$ with no stated cluster count remain
unstandardizable by design, rather than having a cluster count guessed for them.
\item \textbf{Our analyses are far more homogeneous than the humans'.} Within a
condition the standard-error coefficient of variation is $0.002$--$0.09$ against
the crowd's $0.58$--$0.65$, and roughly a third of properly-clustered human
analyses report $|t| < 4$ against about 2\% of ours. No sampler or
standardization change reaches this. It is a property of the analyst, and it is
the central finding rather than a defect.
\end{itemize}

\bibliographystyle{plain}
\bibliography{refs}

\end{document}
;
% also compiles standalone (the guard falls back to article when the workshop
% style is not on TEXINPUTS).
\documentclass{article}

\IfFileExists{neurips_2026.sty}{\usepackage[dblblindworkshop]{neurips_2026}}{%
  % standalone fallback: lmodern gives microtype a scalable font to expand
  \usepackage[margin=1in]{geometry}\usepackage{lmodern}}

\usepackage[utf8]{inputenc}
\usepackage[T1]{fontenc}
\usepackage{hyperref}
\usepackage{url}
\usepackage{booktabs}
\usepackage{amsfonts}
\usepackage{amsmath}
\usepackage{amssymb}
\usepackage{microtype}
\usepackage{graphicx}
\usepackage{xcolor}

\title{Branchpoint: Towards Structured Multiverse Analysis}

\author{%
  Anonymous Author(s) \\
  Affiliation \\
  \texttt{email} \\
}

\begin{document}
\maketitle

\section{Experiment setup}
\label{sec:setup}

This section covers everything outside the sampler that the results depend on:
the datasets, what the analyst is asked, how the model is called, the arms, and
what happens to a result once it exists. The sampler itself is described in
Section~\ref{sec:setup} of the main text; it detects contestable decisions in a
recorded analysis, forks them under a neutral intervention, and grows a ban
list. It makes no validity judgement and never sees a standardization rule.

A third dataset (MARP, religiosity and well-being \cite{hoogeveen2023many})
remains configured but is \emph{dropped throughout}: not run, not standardized,
not plotted, not reported. Nothing below applies to it.

\subsection{Datasets}
\label{sec:setup-data}

Two ecology datasets, each with a published many-analyst comparison, both from
ManyEcoEvo \cite{gould2025same}.

\begin{table}[h]
\centering
\small
\begin{tabular}{lll}
\toprule
 & Blue tit (\texttt{bt}) & Eucalyptus (\texttt{eu}) \\
\midrule
Dimensions            & $3{,}720 \times 42$ & $351 \times 43$ \\
Observation unit      & one chick & one quadrat in one round \\
Clustering            & chicks within ${\sim}330$ rearing nests & quadrats within 18 properties \\
Human analyses coded  & 174 & 128 \\
Usable for $Z_r$      & 111 & 85 \\
\bottomrule
\end{tabular}
\caption{The two datasets. Blue-tit chicks are nested within rearing broods;
Eucalyptus quadrats are not nested in any comparable way, and physical quadrats
do \emph{not} recur across survey rounds. That difference propagates all the way
to the degrees-of-freedom rule in Section~\ref{sec:setup-layer2}.}
\label{tab:datasets}
\end{table}

Analysts answer the published question verbatim: \emph{``To what extent is the
growth of nestling blue tits (Cyanistes caeruleus) influenced by competition
with siblings?''} and \emph{``How does grass cover influence Eucalyptus spp.\
seedling recruitment?''}

\paragraph{Two files per dataset, doing different jobs.}
The analyst-facing data is what the model analyses. A separate comparator file
carries each human team's reported coefficient, standard error and adjusted
degrees of freedom, from which the crowd distribution is derived. Conflating
them is a live hazard, and the two are never interchangeable.

\paragraph{The Eucalyptus design description matters.}
An earlier snapshot of the question stated that fixed quadrats recur across
survey rounds and that the file has 43 columns. Neither is true: quadrat labels
are reused and are not persistent identifiers, and the supplied file has 38
usable columns. Analysts given the false description fit quadrat-clustered
repeated-measures models, which report an asymptotic Wald $z$ with no finite
error df and are then unstandardizable. Every run reported here uses the
corrected snapshot, hash-recorded per condition alongside the results.

\subsection{What the analyst is told}
\label{sec:setup-prompt}

One prompt per dataset, with a fixed structure: the published research question;
a summary of the study design; a column guide; the declared estimand in plain
words; and a frame policy stating what counts as in-frame. Then the working
instruction: analyse the data, report the fitted analysis's focal estimate and
uncertainty \emph{in its native model scale}, work in small checkable Python
scripts, and report a note describing exactly what model, variables and
exclusions were used.

\paragraph{The analyst is never told the standardization rule.}
No Fisher-$z$ formula, no $t$/df transformation map, no degrees-of-freedom
convention. This boundary is marked in code and structurally tested. It matters
because the transformation is exactly the kind of downstream convention that
would contaminate the analytic choices being measured.

This mirrors the benchmark. ManyEcoEvo analysts received the one-sentence
question and the data and nothing else; the information needed for
standardization was elicited \emph{afterwards}, in a separate follow-up survey
with published per-dataset instructions \cite{gould2025same}. Our pipeline
currently reconstructs that information from the analyst's free-text note rather
than asking for it, which is a limitation we return to in
Section~\ref{sec:setup-limits}.

\subsection{How the model is called}
\label{sec:setup-calls}

\paragraph{Loop.}
A files-as-state analysis loop gives the model two tools: \texttt{run\_code},
which writes and runs one Python analysis step in the project directory, and
\texttt{submit\_report}, which finishes with the focal estimate and standard
error. Every step appends to a durable transcript. The model builds its analysis
across small checkable steps rather than one script.

\paragraph{Provider interface.}
All arms use the provider-native interface. Per call we record a secret-free
request object, the complete SDK-serialized response, opaque reasoning blocks,
usage, and every host tool result. This is an SDK-level record, not an
HTTP-wire capture, and it does not expose hidden model state.

\paragraph{Models.}
Three analysts (\texttt{gpt-5.6-sol}, \texttt{claude-opus-5},
\texttt{glm-5.3}) and one judge (\texttt{gpt-5.6-sol}), all at low reasoning
effort. \texttt{glm-5.3} is pinned to a single provider for continuation so the
native history is stable across a fork. Cost is metered per call, tagged by
role, written to a ledger, and re-derived from token counts at one rate table; a
spend guard aborts a run above a ceiling. The full grid below cost \$105.

\subsection{Arms}
\label{sec:setup-arms}

Three arms, sharing the same loop, tools and question.

\begin{table}[h]
\centering
\small
\begin{tabular}{lp{9.4cm}}
\toprule
Arm & What it does \\
\midrule
Naive & $n$ independent analyses from the same prompt under different seeds. The
control: no decision detection, no forking, no intervention. \\
Branchpoint & Detects decisions in a completed trajectory, copies state after the
complete response, appends a neutral intervention, and continues from the exact
provider-native history. Branch points are drawn at random from the open set. \\
BpDynamic & Identical to Branchpoint except allocation: open decisions are
sorted by branches taken plus strikes, and the least-tried are drawn first, so
budget is spent on uncovered decisions. \\
\bottomrule
\end{tabular}
\caption{The three arms. Branchpoint and BpDynamic differ only in how branch
points are allocated, not in how a branch is produced.}
\label{tab:arms}
\end{table}

The intervention is one neutral sentence naming the decision and listing the
approaches already used: \emph{``For \textlangle decision\textrangle, do not use
any approach in this JSON list: [\ldots]. Choose a different defensible
approach; everything else remains yours to decide.''}

\paragraph{Grid as run.}
$\{\textrm{sol}, \textrm{opus}, \textrm{glm}\} \times \{\textrm{bt},
\textrm{eu}\} \times \{\textrm{Naive}, \textrm{Branchpoint}\}$, plus
$\textrm{sol} \times \{\textrm{bt}, \textrm{eu}\} \times \textrm{BpDynamic}$:
14 conditions, $n = 50$ each, ${\sim}700$ analyses.

\subsection{Layer 2: what happens to a result}
\label{sec:setup-layer2}

Layer 2 is host-side, runs only after Layer 1 is banked, and is the sole owner
of the conversion to a common scale. It never runs or extends an analysis agent.
Four stages, each digest-bound so an unchanged input is never re-judged.

\paragraph{(i) Extract.}
A judge reads what the run printed and returns a fixed schema: number of
observations, the unit, the design, exactly one of a Student $t$ or an
asymptotic Wald $z$ (never relabelled), the residual df, a separate finite
reference df used for inference such as clusters minus one, and any directly
reported standardized effect. The digest binds prompt, schema and judge model.

\paragraph{(ii) Standardize --- deterministic, no model call.}
The partial correlation and its transform are
\begin{equation}
r = \frac{t}{\sqrt{t^{2} + \mathit{df}}},
\qquad
Z_r = \operatorname{atanh}(r),
\qquad
\mathrm{SE}(Z_r) = \frac{1}{\sqrt{\mathit{df}}} .
\label{eq:standardize}
\end{equation}
A row reporting both $t$ and $z$ is blocked as ambiguous. A row with a statistic
but no usable df is blocked, but not immediately dropped: it falls through to a
directly reported effect, taken as-is when already on the Fisher-$z$ scale and
via $\operatorname{atanh}$ when $|{\cdot}| < 1$. Only when no usable direct
effect exists does the blocked reason become the row's terminal state, and the
row carries that exact reason rather than disappearing.

Note that $\mathrm{SE}(Z_r)$ depends only on the degrees of freedom: the
analyst's own reported standard error is deliberately \emph{not} used, following
the same convention as the benchmark \cite{gould2025same}. A consequence worth
stating plainly is that our intervals carry no information beyond df, so within
a condition they are nearly identical.

\paragraph{Degrees of freedom are per-dataset policy.}
Equation~\ref{eq:standardize} requires the df that the comparator itself used,
and the two datasets do not agree. Measured from each dataset's own human
panel, the median ratio of adjusted df to sample size is $0.126$ for blue tit
(90 of 96 analyses on a cluster scale) and $0.936$ for Eucalyptus (40 of 55 on a
row scale). They differ because chicks nest within broods and quadrats do not.

Policy therefore decides \emph{precedence}, not merely a fallback: on a
row-scale comparator the residual df wins, then $n - 2$, then any reference df;
on a cluster-scale comparator the reference df wins, then the residual df. This
matters because $\mathrm{SE} = 1/\sqrt{\mathit{df}}$, so a row on a cluster df
sitting beside a row on a row df receives an interval several-fold different for
the same analysis type, and the \emph{gap between the two scales} is then read
as spread. Before precedence was enforced, Eucalyptus widths were inflated
roughly threefold with no change in the number of rows retained.

Where the comparator is row-scale and a generalized linear model reports an
asymptotic Wald $z$ with no classical error df, $n - 2$ is derived. It is the
conservative end, since it maximizes df and so cannot inflate $|r|$, and the
number of parameters is immaterial at these sample sizes: at $n = 350$,
$z = -1.12$, $Z_r$ is $-0.0602$ at $\mathit{df} = 345$ against $-0.0615$ at
$\mathit{df} = 330$. The df \emph{family} is what matters --- a cluster df of 17
would give $-0.262$, over four times larger. Where the comparator is
cluster-scale no cluster count is ever guessed; a row carrying only a row-scale
df stays on it and is reported as such.

\paragraph{(iii) Scope.}
One judgement per standardizable row, against the declared estimand and frame
only. The prompt is explicit that a different covariate set, model family or
estimator is an analytical choice and \emph{not} a scope change, and that
neither is the aggregation level: individuals, clusters and group means are all
reasonable ways to answer one question, and the human panels contain all three.
In this run the screen removed almost nothing.

\paragraph{(iv) Finalize.}
The reportable spread over rows that are both standardizable and in scope. Every
excluded row is retained with its reason, and the counts are reported
separately, so attrition is always visible beside the width --- which matters,
because the two are not independent.

\paragraph{Manual overrides are not used here.}
The pipeline supports a per-condition override file that short-circuits
\emph{before} the deterministic path, supplying a value and method directly.
None of the 14 conditions reported here carries one; every value below is
produced by Equation~\ref{eq:standardize}. This is worth stating because a
condition carrying such a file has rows that no subsequent code fix can reach.

\subsection{Known limitations}
\label{sec:setup-limits}

\begin{itemize}
\item \textbf{Standardization information is reconstructed, not elicited.} The
benchmark asked its analysts for the statistic and df in a follow-up survey; we
parse a free-text note. Analyses whose author wrote ``$z$'' rather than ``$t$
with df'' were being dropped by reporting convention rather than by anything
analytic. The principled fix is a post-analysis elicitation stage.
\item \textbf{Coverage is capped, not converged.} With $n = 50$ against roughly
200 detected decisions, no decision is explored to exhaustion: nothing closes,
and every run stops at the cap. Reported widths are lower bounds from a partial
sample, and coverage belongs beside them.
\item \textbf{Ban lists are per-trajectory.} The same decision axis explored
from a different parent starts with an empty ban list, so distinct branches can
land on identical specifications.
\item \textbf{Blue-tit attrition is uneven.} One condition retains 10 of 50.
Rows reporting a cluster-robust $z$ with no stated cluster count remain
unstandardizable by design, rather than having a cluster count guessed for them.
\item \textbf{Our analyses are far more homogeneous than the humans'.} Within a
condition the standard-error coefficient of variation is $0.002$--$0.09$ against
the crowd's $0.58$--$0.65$, and roughly a third of properly-clustered human
analyses report $|t| < 4$ against about 2\% of ours. No sampler or
standardization change reaches this. It is a property of the analyst, and it is
the central finding rather than a defect.
\end{itemize}

\bibliographystyle{plain}
\bibliography{refs}

\end{document}
;
% also compiles standalone (the guard falls back to article when the workshop
% style is not on TEXINPUTS).
\documentclass{article}

\IfFileExists{neurips_2026.sty}{\usepackage[dblblindworkshop]{neurips_2026}}{%
  % standalone fallback: lmodern gives microtype a scalable font to expand
  \usepackage[margin=1in]{geometry}\usepackage{lmodern}}

\usepackage[utf8]{inputenc}
\usepackage[T1]{fontenc}
\usepackage{hyperref}
\usepackage{url}
\usepackage{booktabs}
\usepackage{amsfonts}
\usepackage{amsmath}
\usepackage{amssymb}
\usepackage{microtype}
\usepackage{graphicx}
\usepackage{xcolor}

\title{Branchpoint: Towards Structured Multiverse Analysis}

\author{%
  Anonymous Author(s) \\
  Affiliation \\
  \texttt{email} \\
}

\begin{document}
\maketitle

\section{Experiment setup}
\label{sec:setup}

This section covers everything outside the sampler that the results depend on:
the datasets, what the analyst is asked, how the model is called, the arms, and
what happens to a result once it exists. The sampler itself is described in
Section~\ref{sec:setup} of the main text; it detects contestable decisions in a
recorded analysis, forks them under a neutral intervention, and grows a ban
list. It makes no validity judgement and never sees a standardization rule.

A third dataset (MARP, religiosity and well-being \cite{hoogeveen2023many})
remains configured but is \emph{dropped throughout}: not run, not standardized,
not plotted, not reported. Nothing below applies to it.

\subsection{Datasets}
\label{sec:setup-data}

Two ecology datasets, each with a published many-analyst comparison, both from
ManyEcoEvo \cite{gould2025same}.

\begin{table}[h]
\centering
\small
\begin{tabular}{lll}
\toprule
 & Blue tit (\texttt{bt}) & Eucalyptus (\texttt{eu}) \\
\midrule
Dimensions            & $3{,}720 \times 42$ & $351 \times 43$ \\
Observation unit      & one chick & one quadrat in one round \\
Clustering            & chicks within ${\sim}330$ rearing nests & quadrats within 18 properties \\
Human analyses coded  & 174 & 128 \\
Usable for $Z_r$      & 111 & 85 \\
\bottomrule
\end{tabular}
\caption{The two datasets. Blue-tit chicks are nested within rearing broods;
Eucalyptus quadrats are not nested in any comparable way, and physical quadrats
do \emph{not} recur across survey rounds. That difference propagates all the way
to the degrees-of-freedom rule in Section~\ref{sec:setup-layer2}.}
\label{tab:datasets}
\end{table}

Analysts answer the published question verbatim: \emph{``To what extent is the
growth of nestling blue tits (Cyanistes caeruleus) influenced by competition
with siblings?''} and \emph{``How does grass cover influence Eucalyptus spp.\
seedling recruitment?''}

\paragraph{Two files per dataset, doing different jobs.}
The analyst-facing data is what the model analyses. A separate comparator file
carries each human team's reported coefficient, standard error and adjusted
degrees of freedom, from which the crowd distribution is derived. Conflating
them is a live hazard, and the two are never interchangeable.

\paragraph{The Eucalyptus design description matters.}
An earlier snapshot of the question stated that fixed quadrats recur across
survey rounds and that the file has 43 columns. Neither is true: quadrat labels
are reused and are not persistent identifiers, and the supplied file has 38
usable columns. Analysts given the false description fit quadrat-clustered
repeated-measures models, which report an asymptotic Wald $z$ with no finite
error df and are then unstandardizable. Every run reported here uses the
corrected snapshot, hash-recorded per condition alongside the results.

\subsection{What the analyst is told}
\label{sec:setup-prompt}

One prompt per dataset, with a fixed structure: the published research question;
a summary of the study design; a column guide; the declared estimand in plain
words; and a frame policy stating what counts as in-frame. Then the working
instruction: analyse the data, report the fitted analysis's focal estimate and
uncertainty \emph{in its native model scale}, work in small checkable Python
scripts, and report a note describing exactly what model, variables and
exclusions were used.

\paragraph{The analyst is never told the standardization rule.}
No Fisher-$z$ formula, no $t$/df transformation map, no degrees-of-freedom
convention. This boundary is marked in code and structurally tested. It matters
because the transformation is exactly the kind of downstream convention that
would contaminate the analytic choices being measured.

This mirrors the benchmark. ManyEcoEvo analysts received the one-sentence
question and the data and nothing else; the information needed for
standardization was elicited \emph{afterwards}, in a separate follow-up survey
with published per-dataset instructions \cite{gould2025same}. Our pipeline
currently reconstructs that information from the analyst's free-text note rather
than asking for it, which is a limitation we return to in
Section~\ref{sec:setup-limits}.

\subsection{How the model is called}
\label{sec:setup-calls}

\paragraph{Loop.}
A files-as-state analysis loop gives the model two tools: \texttt{run\_code},
which writes and runs one Python analysis step in the project directory, and
\texttt{submit\_report}, which finishes with the focal estimate and standard
error. Every step appends to a durable transcript. The model builds its analysis
across small checkable steps rather than one script.

\paragraph{Provider interface.}
All arms use the provider-native interface. Per call we record a secret-free
request object, the complete SDK-serialized response, opaque reasoning blocks,
usage, and every host tool result. This is an SDK-level record, not an
HTTP-wire capture, and it does not expose hidden model state.

\paragraph{Models.}
Three analysts (\texttt{gpt-5.6-sol}, \texttt{claude-opus-5},
\texttt{glm-5.3}) and one judge (\texttt{gpt-5.6-sol}), all at low reasoning
effort. \texttt{glm-5.3} is pinned to a single provider for continuation so the
native history is stable across a fork. Cost is metered per call, tagged by
role, written to a ledger, and re-derived from token counts at one rate table; a
spend guard aborts a run above a ceiling. The full grid below cost \$105.

\subsection{Arms}
\label{sec:setup-arms}

Three arms, sharing the same loop, tools and question.

\begin{table}[h]
\centering
\small
\begin{tabular}{lp{9.4cm}}
\toprule
Arm & What it does \\
\midrule
Naive & $n$ independent analyses from the same prompt under different seeds. The
control: no decision detection, no forking, no intervention. \\
Branchpoint & Detects decisions in a completed trajectory, copies state after the
complete response, appends a neutral intervention, and continues from the exact
provider-native history. Branch points are drawn at random from the open set. \\
BpDynamic & Identical to Branchpoint except allocation: open decisions are
sorted by branches taken plus strikes, and the least-tried are drawn first, so
budget is spent on uncovered decisions. \\
\bottomrule
\end{tabular}
\caption{The three arms. Branchpoint and BpDynamic differ only in how branch
points are allocated, not in how a branch is produced.}
\label{tab:arms}
\end{table}

The intervention is one neutral sentence naming the decision and listing the
approaches already used: \emph{``For \textlangle decision\textrangle, do not use
any approach in this JSON list: [\ldots]. Choose a different defensible
approach; everything else remains yours to decide.''}

\paragraph{Grid as run.}
$\{\textrm{sol}, \textrm{opus}, \textrm{glm}\} \times \{\textrm{bt},
\textrm{eu}\} \times \{\textrm{Naive}, \textrm{Branchpoint}\}$, plus
$\textrm{sol} \times \{\textrm{bt}, \textrm{eu}\} \times \textrm{BpDynamic}$:
14 conditions, $n = 50$ each, ${\sim}700$ analyses.

\subsection{Layer 2: what happens to a result}
\label{sec:setup-layer2}

Layer 2 is host-side, runs only after Layer 1 is banked, and is the sole owner
of the conversion to a common scale. It never runs or extends an analysis agent.
Four stages, each digest-bound so an unchanged input is never re-judged.

\paragraph{(i) Extract.}
A judge reads what the run printed and returns a fixed schema: number of
observations, the unit, the design, exactly one of a Student $t$ or an
asymptotic Wald $z$ (never relabelled), the residual df, a separate finite
reference df used for inference such as clusters minus one, and any directly
reported standardized effect. The digest binds prompt, schema and judge model.

\paragraph{(ii) Standardize --- deterministic, no model call.}
The partial correlation and its transform are
\begin{equation}
r = \frac{t}{\sqrt{t^{2} + \mathit{df}}},
\qquad
Z_r = \operatorname{atanh}(r),
\qquad
\mathrm{SE}(Z_r) = \frac{1}{\sqrt{\mathit{df}}} .
\label{eq:standardize}
\end{equation}
A row reporting both $t$ and $z$ is blocked as ambiguous. A row with a statistic
but no usable df is blocked, but not immediately dropped: it falls through to a
directly reported effect, taken as-is when already on the Fisher-$z$ scale and
via $\operatorname{atanh}$ when $|{\cdot}| < 1$. Only when no usable direct
effect exists does the blocked reason become the row's terminal state, and the
row carries that exact reason rather than disappearing.

Note that $\mathrm{SE}(Z_r)$ depends only on the degrees of freedom: the
analyst's own reported standard error is deliberately \emph{not} used, following
the same convention as the benchmark \cite{gould2025same}. A consequence worth
stating plainly is that our intervals carry no information beyond df, so within
a condition they are nearly identical.

\paragraph{Degrees of freedom are per-dataset policy.}
Equation~\ref{eq:standardize} requires the df that the comparator itself used,
and the two datasets do not agree. Measured from each dataset's own human
panel, the median ratio of adjusted df to sample size is $0.126$ for blue tit
(90 of 96 analyses on a cluster scale) and $0.936$ for Eucalyptus (40 of 55 on a
row scale). They differ because chicks nest within broods and quadrats do not.

Policy therefore decides \emph{precedence}, not merely a fallback: on a
row-scale comparator the residual df wins, then $n - 2$, then any reference df;
on a cluster-scale comparator the reference df wins, then the residual df. This
matters because $\mathrm{SE} = 1/\sqrt{\mathit{df}}$, so a row on a cluster df
sitting beside a row on a row df receives an interval several-fold different for
the same analysis type, and the \emph{gap between the two scales} is then read
as spread. Before precedence was enforced, Eucalyptus widths were inflated
roughly threefold with no change in the number of rows retained.

Where the comparator is row-scale and a generalized linear model reports an
asymptotic Wald $z$ with no classical error df, $n - 2$ is derived. It is the
conservative end, since it maximizes df and so cannot inflate $|r|$, and the
number of parameters is immaterial at these sample sizes: at $n = 350$,
$z = -1.12$, $Z_r$ is $-0.0602$ at $\mathit{df} = 345$ against $-0.0615$ at
$\mathit{df} = 330$. The df \emph{family} is what matters --- a cluster df of 17
would give $-0.262$, over four times larger. Where the comparator is
cluster-scale no cluster count is ever guessed; a row carrying only a row-scale
df stays on it and is reported as such.

\paragraph{(iii) Scope.}
One judgement per standardizable row, against the declared estimand and frame
only. The prompt is explicit that a different covariate set, model family or
estimator is an analytical choice and \emph{not} a scope change, and that
neither is the aggregation level: individuals, clusters and group means are all
reasonable ways to answer one question, and the human panels contain all three.
In this run the screen removed almost nothing.

\paragraph{(iv) Finalize.}
The reportable spread over rows that are both standardizable and in scope. Every
excluded row is retained with its reason, and the counts are reported
separately, so attrition is always visible beside the width --- which matters,
because the two are not independent.

\paragraph{Manual overrides are not used here.}
The pipeline supports a per-condition override file that short-circuits
\emph{before} the deterministic path, supplying a value and method directly.
None of the 14 conditions reported here carries one; every value below is
produced by Equation~\ref{eq:standardize}. This is worth stating because a
condition carrying such a file has rows that no subsequent code fix can reach.

\subsection{Known limitations}
\label{sec:setup-limits}

\begin{itemize}
\item \textbf{Standardization information is reconstructed, not elicited.} The
benchmark asked its analysts for the statistic and df in a follow-up survey; we
parse a free-text note. Analyses whose author wrote ``$z$'' rather than ``$t$
with df'' were being dropped by reporting convention rather than by anything
analytic. The principled fix is a post-analysis elicitation stage.
\item \textbf{Coverage is capped, not converged.} With $n = 50$ against roughly
200 detected decisions, no decision is explored to exhaustion: nothing closes,
and every run stops at the cap. Reported widths are lower bounds from a partial
sample, and coverage belongs beside them.
\item \textbf{Ban lists are per-trajectory.} The same decision axis explored
from a different parent starts with an empty ban list, so distinct branches can
land on identical specifications.
\item \textbf{Blue-tit attrition is uneven.} One condition retains 10 of 50.
Rows reporting a cluster-robust $z$ with no stated cluster count remain
unstandardizable by design, rather than having a cluster count guessed for them.
\item \textbf{Our analyses are far more homogeneous than the humans'.} Within a
condition the standard-error coefficient of variation is $0.002$--$0.09$ against
the crowd's $0.58$--$0.65$, and roughly a third of properly-clustered human
analyses report $|t| < 4$ against about 2\% of ours. No sampler or
standardization change reaches this. It is a property of the analyst, and it is
the central finding rather than a defect.
\end{itemize}

\bibliographystyle{plain}
\bibliography{refs}

\end{document}
